At the heart of every computer lies the Central Processing Unit (CPU),
which is
an electronic circuit that carries out the basic arimethic- and logic
calculation as well as process and redirect input and output to other devices
in the computer.

Most modern CPUs are contained on a very small, yet packed intergrated circuit
chip, which can also house memory caches, multiple cores, and other processing
units.

The functionality of all processors is fundamentally the same. The processor
executes some primitive operation by fetching an instruction in the form of
binary signals, act upon the instruction and store the result in either
one of the its registers, or in the main memory.

Because it is hard and time-consuming to write programs using only the assembly language (writing
only the instructions), compilers are used to remove this complexity. A
compiler takes a high-level language, such as C, C++ or Java, and creates the
corresponding assembly language program for the specific architecture,
containing the instructions. This assembly language file is in turn assembled
or translated to binary, that the particular CPU can understand.

Besides hiding the complexity of the underlying architecture away from the
programmer, it can usually also compile programs to multiple architectures as
well as optimising the code to run faster.

\subsection{Instruction Set Architectures}
The instructions supported by a particular processor is determined by the
Instruction Set Architecture (ISA), which is the specification of how the CPU
works. An ISA determines the instructions
supported, the registers available, memory architecture, addressing modes as
well as handling of interrupts.

There exists many different types of ISAs, with both their advantages and
disadvantages. For example, some architectures have a very few instructions
and
registers, which is very practical for small embedded devices, whereas large
servers might make use of a large array of registers for complex computations.

Besides the current use of an architecture, designers must also take into
account its future uses and applications. As the world of computation is ever
growing and evolving at exponential rates, the architectures must be up to the
challenge of future computing. Introducing a completely new architecture to
the market is very troublesome, and causes a list of problems.
One of the main issues is that old software written for older architectures
will no longer work, and it requires to be either recompiled, rewritten,
or even emulated. One such example is the Intel Itanium (IA-64) architecture,
which had a very bad marked reception due to its lack of backwards compatibility
with the x86 architecture. The emulation of the architecture on IA-64 yielded
suboptimal performance and ultimately lost to the AMD x86\_64, which in turn was
compatible.\cite{anandtech:1854}

Indeed, there are a lot of factors to take into account when designing a new
architecture, and every decision has big implications on the future of the
whole ISA.

\subsection{MIPS}
MIPS is a reduced instruction set architecture (RISC) instruction set
architecture, developed by MIPS technologies in 1981. As an RISC
architecture, MIPS aims to implement only the most essential instructions, so
that they in return can get highly optimised. This is the based on the RISC
philosophy, that by implementing only the most common instructions, the
architects can simplify the design and speed up the crucial parts of the
instructions. This enables the processor to execute programs faster, but also
removes a lot of complexity of implementing large programs.

In contrast to RISC, complex instruction set architecture (CISC) aims
to reduce the number of instructions needed to execute a program by
implementing instructions packed with functionality. This means that a single
instruction in CISC can execute several operations at once, such as loading from memory, arithmetic
and storing. While complex programs indeed execute faster on a CISC architecture,
the burden of implementing efficient and maintainable code and compilers can
outweigh its advantages. \cite{Patterson:1980:CRI:641914.641917}

The first versions of MIPS were 32-bit, moving on to creating 64-bit versions
for computationally heavy applications. It has 31 general purpose (GP)
registers, with additional registers per co-processing unit. MIPS supports up
to 4 coprocessors (COP), with each COP extending on the basic functionality. The
only required co-processor is COP0, which is the System Control Coprocessor.
Other co-processors usually available, but not required, are the floating-point operation
coprocessor COP1 (FPU), the user defined COP2.


\subsection{MIPS Instruction Set}
Each instruction in MIPS is 32-bit long, aligned to word. This simplifies the
instruction fetching, decoding, as well as disassembly of the program.

In MIPS, the instructions have 3 basic formats
:
\begin{figure}[ht]
\centering
\subfigure[R-Format]{%
	\includegraphics{cpu_architecture/r_format.eps}
	\label{fig:instruction_r_format}}
\subfigure[I-Format]{%
	\includegraphics{cpu_architecture/i_format.eps}
	\label{fig:instruction_i_format}}
\subfigure[J-Format]{%
	\includegraphics{cpu_architecture/j_format.eps}
	\label{fig:instruction_j_format}}
\label{fig:instruction_formats}
\end{figure}

It is clear that as they have common fields, mainly the opcode field, they are
easily distinguishable.

